\documentclass[a4paper,10pt]{report}
\usepackage[a4paper,width=150mm,top=25mm,bottom=25mm]{geometry}

\usepackage[utf8]{inputenc}
\usepackage{xspace}
\usepackage[portuguese]{babel}

% o comando say
\usepackage{dirtytalk}

%para que o primeiro paragrafo indente
\usepackage{indentfirst}

%Fontes especias de matemática
\usepackage{amsfonts}

%hide hilelinks esconde os links de terem bordas feias
\usepackage[hidelinks]{hyperref}
% pacote para os headers e footers fixes
\usepackage{fancyhdr}

%para poder ter os compactitem lists
\usepackage{paralist}

%ter acesso a imagens
\usepackage{graphicx}
\graphicspath{{images/}}

\usepackage{setspace}
\usepackage{afterpage}
% CORES

%tratar das cores
\usepackage{xcolor}
\usepackage{pagecolor,lipsum}% http://ctan.org/pkg/{pagecolor,lipsum}

% definição de uma cor
% \definecolor{name}{model}{color-spec}
\definecolor{aquamarine}{rgb}{0, 181, 190}




% meter o fundo em pretos e os caracters em branco
% \pagecolor{black}
% \color{white}
% \newcommand{\linkcolor}{aquamarine}
%link color for white mode
\newcommand{\linkcolor}{blue}




\usepackage{xcolor}
\usepackage{todonotes}
% \usepackage{biblatex}
% \addbibresource{referencias.bib}

%MACROS
\newcommand{\lk}[2]{\href{#1}{\textcolor{\linkcolor}{#2}}}
\newcommand{\chapterf}[1]{\chapter{#1}\thispagestyle{fancy}}


% Informações do autor
\newcommand{\myauthor}{João Pedro Neves Gonçalves}
\newcommand{\mynumerodealuno}{2018014306}

%valores da palestra
\newcommand{\palestratitulo} {O Colégio de Engenharia Informática:\\
                                oportunidades, domínios de intervenção e competências}
\newcommand{\palestraautor} {Isabel Pedrosa \& Jorge Sá Silva}
\newcommand{\palestradata} {7 de Abril de 2021}
\newcommand{\palestranumero} {4}


% set dos headers e footers
\pagestyle{fancy}
\fancyhf{}
\lhead{\mytitle{}}
\rhead {\thechapter}
\rfoot{\thepage}


%set do espaçamento entre linhas
\setstretch{1.25}
%set indentação
\setlength\parindent{0.75cm}

%set font family to arial
\usepackage{helvet}
\renewcommand{\familydefault}{\sfdefault}

\newcommand \mytitle {Relatório da Palestra nº \palestranumero}

\title{\mytitle}

\author{\myauthor}

\newcommand\myemptypage{
    \newpage
    \null
    \thispagestyle{empty}
    % \addtocounter{page}{-1}
    \newpage
    }


\begin{document}


\begin{titlepage}
    \begin{flushright}
        \textbf{Departamento de Engenharia Informática e de Sistemas
            Instituto Superior de Engenharia de Coimbra}

        Instituto Politécnico de Coimbra
    \end{flushright}
    \begin{center}
        \vspace{0.5cm}
        \begin{large}
            \newcommand{\spacetitle}{0.0cm}

            Curso Licenciatura Engenharia Informática - Diurno
            \vspace{\spacetitle}

            Ramo Desenvolvimento de Aplicações
            \vspace{\spacetitle}

            Unidade Curricular de Ética e Deontologia
            \vspace{\spacetitle}

            Ano Letivo de 2020/2021
            \vspace{\spacetitle}

            \vfill

            Palestra nº \palestranumero

            \vspace{\spacetitle}
            \vspace{\spacetitle}
            \textbf{\palestratitulo}
            \vspace{\spacetitle}
            \vspace{\spacetitle}
            
            Realizado por \palestraautor

            Realizado em \palestradata

            \vfill

            \includegraphics[width=0.4\textwidth]{ISEC}
            \vfill

            \myauthor

            Número de Aluno: \mynumerodealuno

            Coimbra, \today
            \vspace{\spacetitle}

            \myauthor

        \end{large}
    \end{center}
\end{titlepage}

\myemptypage{}

\tableofcontents{}
\thispagestyle{empty}

\newpage
\myemptypage{}
\newpage

\pagenumbering{arabic}
\section{Resumo}
Este trabalho é feito no âmbito de Ética e Deontologia, em que se fez um relatório sobre ``O Colégio de Engenharia Informática: oportunidades, domínios de intervenção e competências'' que foi a \palestranumero ª palestra apresentada na disciplina.

Está palestra foi apresentada pelo Engenheiro Jorge Sá Silva  e pela Engenheira Isabel Pedrosa em que se falou sobre a Ordem dos Engenheiros, o porque de fazer parte desta, eventos em que está relacionada, foi também apresentado um foco especial no colégio de informática e também foi apresentado alguns algumas opiniões sobre a Ordem e relacionados no final.


\myemptypage{}

\chapterf{Introdução}

Engenharia trata de vários pontos na sociedade e é preciso de uma ordem para tratar de proteger os direitos e deveres dos elementos que a praticam, neste documento tem se relatado a palestra que a apresenta.


\myemptypage{}
\chapterf{Descrição da Palestra}


%analise

\section{Apresentação}
O engenheiro Jorge Sá Silva começa a palestra por apresentar varias questões que as pessoas fazem quando é para se inscrever na ordem dos engenheiros, estas são:

\begin{enumerate}
    \item Porquê inscrever-me na Ordem dos engenheiros?
    \item Porquê pagar os 10€ por mês?
    \item Porquê inscrever-me se a empresa que me vai contratar não exige isso?
    \item Acabo o curso e ainda tenho de fazer um estágio de vários meses sem ganhar dinheiro só para entrar para a OE?
    \item Devo inscrever-me na Ordem dos engenheiros Técnicos, sendo eu aluno de um Instituto Politécnico?
\end{enumerate}

E refere o que irá apresentar sobre a ordem:

\begin{enumerate}
    \item História
    \item Objetivos
    \item Organização e os seus membros
    \item O que fazem
    \item Benefícios
    \item Outros
\end{enumerate}

\newpage
\section{Historia}
\begin{itemize}
    \item Originou da Associação dos engenheiros Civis Portugueses
    \item Fundada em 1969
    \item O titulo de engenheiro foi definido pela primeira vez em 26 de Julho de 1926
    \item Foi formalizada em 1936
    \item A ordem é atualmente suportada por 12 colégios, em varias regiões do pais
\end{itemize}

\section{O que fazem}

O principal objetivo da ordem é contribuir com o progresso à Engenharia, tem a responsabilidade de atribuir o titulo oficial de engenheiro, é também está que sela os direitos e deveres dos seus membros.

Ela participa no ensino e desenvolvimento da engenheira como assistir no desenvolvimento da estrutura das cadeiras para os engenheiros, Desenvolvimento com outras instituições e apoiar projetos de éticos.

\section{A Organização}

A organização tem 9  órgãos nacionais e 5 regionais com vários conselhos e uma assembleia de representantes.

\section{Colégios}

\begin{itemize}
    \item Agronómica
    \item Ambiente
    \item Civil
    \item Eletrotécnica
    \item Florestal
    \item Geográfica
    \item Geológia e de Minas
    \item Informática
    \item Materiais
    \item Mecânica
    \item Naval
\end{itemize}

Estes podem ser vistos com mais detalhe 
\lk{https://www.ordemengenheiros.pt/pt/a-ordem/colegios-e-especialidades/}{neste site} 
\\
(https://www.ordemengenheiros.pt/pt/a-ordem/colegios-e-especialidades/).



\section{Membros}

Existem 2 tipos de membros:

\begin{itemize}
    \item Membros Estudantes
    \item Membros estagiários e efetivos
    \begin{itemize}
        \item Membros Efetivos
        \item Membros Seniores
        \item Membros Conselheiros
        \item Membros Honorários
    \end{itemize}
\end{itemize}

\section{Respondendo as questões iniciais}

Resumindo, a razão para entrar na ordem é porque esta instituição tem o cuidado de defender a profissão dos engenheiros, é uma ordem profissional/prestígio e é onde a sociedade obtém aconselhamento e assistência dentro da área. 

\section{Processo de admissão e qualificação}

A admissão é feita através da seguintes formas:

\begin{itemize}
    \item Se for do grau mestre pode fazer um estágio de 6 meses.
    \item Se for do grau licenciado pode fazer um estagio de 18 meses.
    \item Também pode fazer um curso de ética e deontologia profissional promovido pela Ordem dos engenheiros.
\end{itemize}

Notas:

\begin{itemize}
    \item Se já tiver experiência na área de 5 ou 6 anos poderá fazer dispensa do estagio.
    \item Normalmente os estágios são remunerados.
\end{itemize} 


\section{Níveis de qualificação}
\begin{itemize}
    \item Engenheira de nível 1 (licenciatura em engenharia)
    \item Engenheira de nível 2 (mestrado em engenharia)
\end{itemize}


\section{Prestigio/Benefícios}
\begin{itemize}
    \item Formação
    \item Seguro da responsabilidade Civil
    \item Informação Jurídica
    \item Revista da ordem dos engenheiros
    \item Bolsa de emprego
    \item Protocolos com diversas instituições e empresas que dão benefícios aos engenheiros
\end{itemize}

\section{Regalias para os membros}

\begin{itemize}
    \item Existem descontos em varias empresas.
    \item Eventos sociais
        \begin{itemize}
            \item A mulher na engenharia
            \begin{itemize}
                \item As Mulheres na Ciência
                \item Girls on IT
            \end{itemize}
            \item O melhor estágio
        \end{itemize}
\end{itemize}

\section{Colégio de Informática}
\begin{itemize}
    \item Um dos colégios
    \item 20.000 Engenheiros
\end{itemize}





\subsection{Atos de engenharia Informática}

Está dividida em 6 áreas:

\begin{itemize}
    \item Análise e validação de requisitos de soluções
    \item Programação de soluções
    \item Informáticas
    \item Integração de soluções
    \item Análise e validação de requisitos de infraestruturas de computação, comunicações e serviços
    \item Definição e modelação de arquiteturas de infraestruturas de computação, comunicações e serviços
    \item Definição e documentação de planos de gestão de projetos de sistemas de Informação
    \item Implementação de planos de gestão e auditoria de níveis de serviços, qualidade, risco e segurança em sistemas de informação
\end{itemize}

\section{Engenheiros pelo mundo}

Uma app implementada pelo colégio informático da zona centro, feita para descobrir os engenheiros espalhados pelo mundo.



%critica
\chapterf{Análise Crítica}

\section{Apresentação}

Eu gostei muito da apresentação que o engenheiro Jorge fez, foi muito informativa e convivente para entrar na ordem, resumiu bem os assuntos e apresentou a um ótimo ritmo.



\myemptypage{}
\rhead{\thechapter}
\chapterf{Conclusão}


Foi uma ótima palestra sobre a ordem em que honestamente convenceu me a querer entrar nela pois para mim ter o apoio Jurídico e outros benefícios parecem ser bastante positivos pelos 10€ mensais.


\end{document}
